\begin{dang}{Lí thuyết các loại tia}
\Anlythuyet{
\begin{itemize}
\item Do có tần số và bước sóng khác nhau nên sóng điện từ có những tính chất rất khác nhau: có thể nhìn thấy hoặc không nhìn thấy, có khả năng đâm xuyên khác nhau, cách phát khác nhau. 
\item Các tia có bước sóng càng ngắn (tia X, tia gamma) có tính đâm xuyên càng mạnh, dễ tác dụng lên kính ảnh, dễ làm phát quang các chất và dễ ion hóa không khí. Trong khi đó, với các tia có bước sóng dài thì ta dễ dàng quan sát hiện tượng giao thoa.
\end{itemize}
}
\end{dang}